% Options for packages loaded elsewhere
\PassOptionsToPackage{unicode}{hyperref}
\PassOptionsToPackage{hyphens}{url}
\PassOptionsToPackage{dvipsnames,svgnames,x11names}{xcolor}
\documentclass[
  10pt,
  fontset=fandol]{ctexart}
\usepackage{xcolor}
\usepackage[margin=16mm]{geometry}
\usepackage{amsmath,amssymb}
\setcounter{secnumdepth}{-\maxdimen} % remove section numbering
\usepackage{iftex}
\ifPDFTeX
  \usepackage[T1]{fontenc}
  \usepackage[utf8]{inputenc}
  \usepackage{textcomp} % provide euro and other symbols
\else % if luatex or xetex
  \usepackage{unicode-math} % this also loads fontspec
  \defaultfontfeatures{Scale=MatchLowercase}
  \defaultfontfeatures[\rmfamily]{Ligatures=TeX,Scale=1}
\fi
\usepackage{lmodern}
\ifPDFTeX\else
  % xetex/luatex font selection
\fi
% Use upquote if available, for straight quotes in verbatim environments
\IfFileExists{upquote.sty}{\usepackage{upquote}}{}
\IfFileExists{microtype.sty}{% use microtype if available
  \usepackage[]{microtype}
  \UseMicrotypeSet[protrusion]{basicmath} % disable protrusion for tt fonts
}{}
\makeatletter
\@ifundefined{KOMAClassName}{% if non-KOMA class
  \IfFileExists{parskip.sty}{%
    \usepackage{parskip}
  }{% else
    \setlength{\parindent}{0pt}
    \setlength{\parskip}{6pt plus 2pt minus 1pt}}
}{% if KOMA class
  \KOMAoptions{parskip=half}}
\makeatother
\usepackage{longtable,booktabs,array}
\newcounter{none} % for unnumbered tables
\usepackage{calc} % for calculating minipage widths
% Correct order of tables after \paragraph or \subparagraph
\usepackage{etoolbox}
\makeatletter
\patchcmd\longtable{\par}{\if@noskipsec\mbox{}\fi\par}{}{}
\makeatother
% Allow footnotes in longtable head/foot
\IfFileExists{footnotehyper.sty}{\usepackage{footnotehyper}}{\usepackage{footnote}}
\makesavenoteenv{longtable}
\setlength{\emergencystretch}{3em} % prevent overfull lines
\providecommand{\tightlist}{%
  \setlength{\itemsep}{0pt}\setlength{\parskip}{0pt}}
\usepackage{amsmath,amssymb,mathtools,needspace}
\xeCJKDeclareCharClass{CJK}{"2032,"0400->"04FF,"2160->"216B,"2460->"2473,"25A0,"25C7,"25CB,"25CF}
\IfFontExistsTF{Arial Unicode MS}{\xeCJKsetup{AutoFallBack=true}\setCJKfallbackfamilyfont{\CJKrmdefault}{Arial Unicode MS}}{}
\setcounter{secnumdepth}{0}
\setlength{\emergencystretch}{2em}
\allowdisplaybreaks
\usepackage{bookmark}
\IfFileExists{xurl.sty}{\usepackage{xurl}}{} % add URL line breaks if available
\urlstyle{same}
\hypersetup{
  pdftitle={Runge-Kutta 方法},
  colorlinks=true,
  linkcolor={Maroon},
  filecolor={Maroon},
  citecolor={Blue},
  urlcolor={Blue},
  pdfcreator={LaTeX via pandoc}}

\title{Runge-Kutta 方法}
\author{}
\date{}

\begin{document}
\maketitle

\subsection{5.3 Runge-Kutta 方法}\label{runge-kutta-ux65b9ux6cd5}

本节介绍一类高精度的单步法------Runge-Kutta 方法。这类方法与下述 Taylor 级数法有着紧密的联系。

\subsubsection{5.3.1 Taylor 级数法}\label{taylor-ux7ea7ux6570ux6cd5}

设初值问题 (5.1.1)、(5.1.2) 有解 \(y=y(x)\)，用 Taylor 展开，有

\[
y(x_{n+1})=y(x_n)+hy'(x_n)+\frac{h^2}{2!}y''(x_n)+\frac{h^3}{3!}y'''(x_n)+\cdots,\tag{5.3.1}
\]

其中 \(y(x)\) 的各阶导数依据所给方程 (5.1.1) 可以用函数 \(f\) 来表达。下面引进函数序列 \(f^{(j)}(x,y)\) 来描述求导过程，即

\[
y'=f\equiv f^{(0)},\qquad
y''=\frac{\partial f^{(0)}}{\partial x}+f\frac{\partial f^{(0)}}{\partial y}\equiv f^{(1)},\qquad
y'''=\frac{\partial f^{(1)}}{\partial x}+f\frac{\partial f^{(1)}}{\partial y}\equiv f^{(2)},\qquad\cdots.
\]

一般地

\[
y^{(j)}=\frac{\partial f^{(j-2)}}{\partial x}+f\frac{\partial f^{(j-2)}}{\partial y}\equiv f^{(j-1)},\qquad j=2,3,\cdots,
\]

而具体写出则有

\[
\begin{cases}
y'=f,\\
y''=\dfrac{\partial f}{\partial x}+f\dfrac{\partial f}{\partial y},\\
y'''=\dfrac{\partial^2f}{\partial x^2}+2f\dfrac{\partial^2f}{\partial x\partial y}
+f^2\dfrac{\partial^2f}{\partial y^2}
+\dfrac{\partial f}{\partial y}\left(\dfrac{\partial f}{\partial x}+f\dfrac{\partial f}{\partial y}\right),\\
\qquad\vdots
\end{cases}\tag{5.3.2}
\]

在展开式 (5.3.1) 的右端截取若干项，并且在 \((x_n,y_n)\) 按式 (5.3.2) 计算系数 \(y^{(j)}(x_n)\) 的近似值 \(y_n^{(j)}\)，结果导出下列 \textbf{Taylor 格式}

\[
y_{n+1}=y_n+hy'_n+\frac{h^2}{2!}y''_n+\cdots+\frac{h^p}{p!}y_n^{(p)},\tag{5.3.3}
\]

其中一阶 Taylor 格式 \((p=1)\)

\[
y_{n+1}=y_n+hy'_n
\]

就是 Euler 格式，而提高 Taylor 格式的阶 \(p\) 即可提高计算结果的精度。显然，\(p\) 阶

\href{../../source-images/pdf-127.jpeg}{原页}

Taylor 格式 (5.3.3) 的局部截断误差为

\[
y(x_{n+1})-y_{n+1}=\frac{h^{p+1}}{(p+1)!}y^{(p+1)}(\xi),\qquad x_n<\xi<x_{n+1},
\]

因此，它按下述定义具有 \(p\) 阶精度。

\textbf{定义 5.1} 如果一种方法的局部截断误差为 \(O(h^{p+1})\)，则称该方法具有 \(p\) 阶精度。

\textbf{例 5.3} 用 Taylor 级数法求解例 5.1 的初值问题 (5.2.2)。

\textbf{解} 直接求导知

\[
\begin{aligned}
y'&=y-\frac{2x}{y},\\
y''&=y'-\frac{2}{y^2}(y-xy'),\\
y'''&=y''+\frac{2}{y^2}(xy''+2y')-\frac{4x(y')^2}{y^3},\\
y^{(4)}&=y'''+\frac{2}{y^2}(xy'''+3y'')-\frac{12y'}{y^3}(xy''+y')+\frac{12x(y')^3}{y^4}.
\end{aligned}
\]

据此用四阶 Taylor 格式即可获得令人满意的结果（仍取步长 \(=0.1\) 计算，结果如表 5.3 所示）。

\textbf{表 5.3}

{\def\LTcaptype{none} % do not increment counter
\begin{longtable}[]{@{}lll@{}}
\toprule\noalign{}
\(x_n\) & \(y_n\) & \(y(x_n)\) \\
\midrule\noalign{}
\endhead
\bottomrule\noalign{}
\endlastfoot
0.1 & 1.0954 & 1.0954 \\
0.2 & 1.1832 & 1.1832 \\
0.3 & 1.2649 & 1.2649 \\
\end{longtable}
}

应当指出，Taylor 格式 (5.3.3) 从形式上看似简单，但具体构造这种格式往往是相当困难的，因为它需要按式 (5.3.2) 提供导数值 \(y_n^{(j)}\)。当阶数提高时，求导过程式 (5.3.2) 可能很复杂，因此 Taylor 级数法通常不直接使用，但可以用它来启发思路。

\subsubsection{5.3.2 Runge-Kutta 方法的基本思想}\label{runge-kutta-ux65b9ux6cd5ux7684ux57faux672cux601dux60f3}

Runge-Kutta 方法实质上是间接地使用 Taylor 级数法的一种方法。

考察差商 \(\dfrac{y(x_{n+1})-y(x_n)}{h}\)，根据微分中值定理，存在 \(0<\theta<1\)，使得

\[
\frac{y(x_{n+1})-y(x_n)}{h}=y'(x_n+\theta h),
\]

于是，利用所给方程 \(y'=f(x,y)\) 得到

\[
y(x_{n+1})=y(x_n)+hf(x_n+\theta h,y(x_n+\theta h)).\tag{5.3.4}
\]

设 \(K^*=f(x_n+\theta h,y(x_n+\theta h))\)，称 \(K^*\) 为区间 \([x_n,x_{n+1}]\) 上的\textbf{平均斜率}。由此可见，只要对平均斜率提供一种算法，那么由式 (5.3.4) 便相应地导出一种计算格式。

Euler 格式 (5.2.1) 简单地取点 \(x_n\) 的斜率值 \(K_1=f(x_n,y_n)\) 作为平均斜率 \(K^*\)，

\href{../../source-images/pdf-128.jpeg}{原页}

精度自然很低。

再考察改进的 Euler 格式 (5.2.10)、(5.2.11)，它可改写成下列平均化的形式：

\[
\begin{cases}
y_{n+1}=y_n+\dfrac{h}{2}(K_1+K_2),\\
K_1=f(x_n,y_n),\\
K_2=f(x_{n+1},y_n+hK_1).
\end{cases}
\]

可见，改进的 Euler 格式可以这样理解：它用 \(x_n\) 与 \(x_{n+1}\) 两个点的斜率值 \(K_1\) 与 \(K_2\) 取算术平均作为平均斜率 \(K^*\)，而 \(x_{n+1}\) 处的斜率值 \(K_2\) 则通过已知信息 \(y_n\) 来预测。

这个处理过程启示我们，如果设法在 \((x_n,x_{n+1})\) 内多预测几个点的斜率值，然后将它们加权平均作为平均斜率 \(K^*\)，则有可能构造出具有更高精度的计算格式。这就是 Runge-Kutta 方法的基本思想。

\subsubsection{5.3.3 二阶 Runge-Kutta 方法}\label{ux4e8cux9636-runge-kutta-ux65b9ux6cd5}

首先推广改进的 Euler 方法，随意考察区间 \((x_n,x_{n+1})\) 内一点

\[
x_{n+p}=x_n+ph,\qquad 0<p\leqslant1,
\]

希望用 \(x_n\) 和 \(x_{n+p}\) 两个点的斜率值 \(K_1\) 和 \(K_2\) 线性组合得到平均斜率 \(K^*\)，即令

\[
y_{n+1}=y_n+h(\lambda_1K_1+\lambda_2K_2),
\]

其中 \(\lambda_1,\lambda_2\) 为待定系数。同改进的 Euler 格式一样，这里仍取 \(K_1=f(x_n,y_n)\)，问题在于该怎样预测 \(x_{n+p}\) 处的斜率值 \(K_2\)。

仿照改进的 Euler 格式，先用 Euler 格式提供 \(y(x_{n+p})\) 的预测值，即

\[
y_{n+p}=y_n+phK_1,
\]

然后再用预测值 \(y_{n+p}\) 通过计算 \(f\) 产生斜率值

\[
K_2=f(x_{n+p},y_{n+p}).
\]

这样设计出的计算格式具有如下形式：

\[
\begin{cases}
y_{n+1}=y_n+h(\lambda_1K_1+\lambda_2K_2),\\
K_1=f(x_n,y_n),\\
K_2=f(x_{n+p},y_n+phK_1).
\end{cases}\tag{5.3.5}
\]

式 (5.3.5) 含有三个待定系数 \(\lambda_1,\lambda_2\) 和 \(p\)，希望适当选取这些系数的值，使得格式具有二阶精度。

根据式 (5.3.3)，二阶 Taylor 格式为

\[
y_{n+1}=y_n+hy'_n+\frac{h^2}{2}y''_n,
\]

或表示为

\[
y_{n+1}=y_n+hf_n+\frac{h^2}{2}(f_x+f\cdot f_y)_n,
\]

其中 \(f_n\) 和 \((f_x+f\cdot f_y)_n\) 的下标 \(n\) 均表示在 \((x_n,y_n)\) 取值。另一方面，有

\href{../../source-images/pdf-129.jpeg}{原页}

\[
K_1=f_n,\qquad K_2=f(x_{n+p},y_n+phK_1)=f_n+ph(f_x+f\cdot f_y)_n+\cdots,
\]

代入式 (5.3.5)，得

\[
y_{n+1}=y_n+(\lambda_1+\lambda_2)hf_n+\lambda_2ph^2(f_x+f\cdot f_y)_n+\cdots.
\]

由此可见，欲使格式 (5.3.5) 具有二阶精度，只要成立

\[
\begin{cases}
\lambda_1+\lambda_2=1,\\
\lambda_2p=\dfrac{1}{2}.
\end{cases}\tag{5.3.6}
\]

满足条件 (5.3.6) 的一族格式 (5.3.5) 统称为\textbf{二阶 Runge-Kutta 格式}。

除了改进的 Euler 格式外，另一种特殊的二阶 Runge-Kutta 格式是所谓\textbf{变形的 Euler 格式}，其形式如下：

\[
\begin{cases}
y_{n+1}=y_n+hK_2,\\
K_1=f(x_n,y_n),\\
K_2=f\left(x_n+\dfrac{h}{2},y_n+\dfrac{h}{2}K_1\right).
\end{cases}
\]

表面上看，变形的 Euler 格式 \(y_{n+1}=y_n+hK_2\) 仅显含一个斜率值 \(K_2\)，但 \(K_2\) 是通过 \(K_1\) 计算出来的，因此每完成一步仍然需要两次计算函数 \(f\) 的值，工作量和改进的 Euler 格式相同。

总之，二阶 Runge-Kutta 方法用多算一次函数值 \(f\) 的办法，避开了二阶 Taylor 级数法所要求的 \(f\) 导数值的计算。在这种意义上可以说，Runge-Kutta 方法实质上是 Taylor 方法的变形。

\subsubsection{5.3.4 三阶 Runge-Kutta 方法}\label{ux4e09ux9636-runge-kutta-ux65b9ux6cd5}

为了进一步提高精度，设除 \(x_n,x_{n+p}\) 外再考察一点 \(x_{n+p}=x_n+qh\)，\(p\leqslant q\leqslant1\)，并用三个点 \(x_n,x_{n+p},x_{n+q}\) 的斜率值 \(K_1,K_2,K_3\) 线性组合得到平均斜率 \(K^*\)，这时计算格式为

\[
y_{n+1}=y_n+h(\lambda_1K_1+\lambda_2K_2+\lambda_3K_3).
\]

其中 \(K_1,K_2\) 仍用格式 (5.3.5) 所取的形式。

为了预测点 \(x_{n+p}\) 处的斜率值 \(K_3\)，在区间 \((x_n,x_{n+q})\) 内有两个斜率值 \(K_1\) 和 \(K_2\) 可以利用。用 \(K_1\) 和 \(K_2\) 线性组合给出区间 \([x_n,x_{n+q}]\) 上的平均斜率，从而得到 \(y(x_{n+q})\) 的预测值 \(y_{n+q}=y_n+qh(rK_1+sK_2)\)，于是，再通过计算函数值 \(f\) 得到

\[
K_3=f(x_{n+q},y_{n+q})=f(x_n+qh,y_n+qh(rK_1+sK_2)).
\]

这样设计出的计算格式具有如下形式：

\[
\begin{cases}
y_{n+1}=y_n+h(\lambda_1K_1+\lambda_2K_2+\lambda_3K_3),\\
K_1=f(x_n,y_n),\\
K_2=f(x_n+ph,y_n+phK_1),\\
K_3=f(x_n+qh,y_n+qh(rK_1+sK_2)).
\end{cases}\tag{5.3.7}
\]

\begin{quote}
\textbf{校注（agent 补充，原书下标错误）}：5.3.4 首段``再考察一点 \(x_{n+p}=x_n+qh\)''及后段``为了预测点 \(x_{n+p}\) 处的斜率值 \(K_3\)''，原图两处均印为 \(n+p\)，已保留。按同页对 \(K_3=f(x_{n+q},y_{n+q})\) 的定义及式 (5.3.7)，这两处所指节点均应为 \(x_{n+q}\)。已放大原印字确认，非 OCR 误识。
\end{quote}

\href{../../source-images/pdf-130.jpeg}{原页}

希望适当选择系数 \(\lambda_1,\lambda_2,\lambda_3\) 和 \(p,q,r,s\)，能使上述格式具有三阶精度。

为便于进行数学演算，引进算子

\[
\mathrm D=\frac{\partial}{\partial x}+f\frac{\partial}{\partial y},\qquad
\mathrm D^2=\frac{\partial^2}{\partial x^2}+2f\frac{\partial^2}{\partial x\partial y}+f^2\frac{\partial^2}{\partial y^2},
\]

则根据式 (5.3.2)，有

\[
y'=f,\qquad y''=\mathrm Df,\qquad y'''=\mathrm D^2f+\frac{\partial f}{\partial y}\mathrm Df,
\]

于是三阶 Taylor 格式

\[
y_{n+1}=y_n+hy'_n+\frac{h^2}{2}y''_n+\frac{h^3}{6}y'''_n
\]

可表示为

\[
y_{n+1}=y_n+hf_n+\frac{h^2}{2}\mathrm Df_n+\frac{h^3}{6}(\mathrm D^2f+f_y\mathrm Df)_n,\tag{5.3.8}
\]

其中 \(f_n,\mathrm Df_n\) 等的下标 \(n\) 均表示在 \((x_n,y_n)\) 取值。

另一方面，再将式 (5.3.7) Taylor 展开，其中

\[
K_1=f_n,
\]

\[
K_2=f(x_n+ph,y_n+phK_1)=f_n+ph\mathrm Df_n+\frac{(ph)^2}{2}\mathrm D^2f_n+\cdots,
\]

\[
K_3=f(x_n+qh,y_n+qh(rK_1+sK_2))=f_n+qh\overline{\mathrm D}f_n+\frac{(qh)^2}{2}\overline{\mathrm D}^{2}f_n+\cdots.
\]

这里，算子 \(\overline{\mathrm D}=\dfrac{\partial}{\partial x}+(rK_1+sK_2)\dfrac{\partial}{\partial y}\)。

若取

\[
r+s=1,\tag{5.3.9}
\]

则易知

\[
\overline{\mathrm D}=\mathrm D+hps\mathrm Df_n\frac{\partial}{\partial y}+\cdots,\tag{5.3.10}
\]

\[
\overline{\mathrm D}^{2}=\mathrm D^2+\cdots,
\]

因此

\[
K_3=f_n+qh\mathrm Df_n+\frac{h^2}{2}(q^2\mathrm D^2f+2pqs\mathrm Df\cdot f_y)_n+\cdots.
\]

将以上 \(K_1,K_2,K_3\) 的展开式一起代入式 (5.3.7)，再令

\[
\lambda_1+\lambda_2+\lambda_3=1,\tag{5.3.11}
\]

则有

\[
\begin{aligned}
y_{n+1}={}&y_n+hf_n+h^2(\lambda_2p+\lambda_3q)\mathrm Df_n\\
&+h^3\left(\frac12\lambda_2p^2\mathrm D^2f+\frac12\lambda_3q^2\mathrm D^2f+\lambda_3pqsf_y\mathrm Df\right)_n+\cdots,
\end{aligned}
\]

于是为了保证它能与三阶 Taylor 格式 (5.3.8) 具有同等精度，只要满足下列方程组：

\[
\begin{cases}
\lambda_2p+\lambda_3q=\dfrac12,\\
\lambda_2p^2+\lambda_3q^2=\dfrac13,\\
\lambda_3pqs=\dfrac16.
\end{cases}\tag{5.3.12}
\]

\href{../../source-images/pdf-131.jpeg}{原页}

满足条件 (5.3.9)、(5.3.11)、(5.3.12) 的一族格式 (5.3.7) 统称\textbf{三阶 Runge-Kutta 格式}。下列 Kutta 格式是其中的一个特例：

\[
\begin{cases}
y_{n+1}=y_n+\dfrac{h}{6}(K_1+4K_2+K_3),\\
K_1=f(x_n,y_n),\\
K_2=f\left(x_n+\dfrac{h}{2},y_n+\dfrac{h}{2}K_1\right),\\
K_3=f(x_n+h,y_n-hK_1+2hK_2).
\end{cases}
\]

\subsubsection{5.3.5 四阶 Runge-Kutta 方法}\label{ux56dbux9636-runge-kutta-ux65b9ux6cd5}

继续上述过程，经过较复杂的数学演算，可以导出各种四阶 Runge-Kutta 格式，下列\textbf{经典格式}是其中常用的一个：

\[
\begin{cases}
y_{n+1}=y_n+\dfrac{h}{6}(K_1+2K_2+2K_3+K_4),\\
K_1=f(x_n,y_n),\\
K_2=f\left(x_n+\dfrac{h}{2},y_n+\dfrac{h}{2}K_1\right),\\
K_3=f\left(x_n+\dfrac{h}{2},y_n+\dfrac{h}{2}K_2\right),\\
K_4=f(x_n+h,y_n+hK_3).
\end{cases}\tag{5.3.13}
\]

四阶 Runge-Kutta 方法的每一步需要四次计算函数值 \(f\)，可以证明其截断误差为 \(O(h^5)\)。其证明极其烦琐，这里从略。

\textbf{例 5.4} 设取步长 \(h=0.2\)，从 \(x=0\) 直到 \(x=1\) 用四阶 Runge-Kutta 方法求解初值问题 (5.2.2)。

\textbf{解} 这里，经典的四阶 Runge-Kutta 格式 (5.3.13) 具有如下形式：

\[
\begin{cases}
y_{n+1}=y_n+\dfrac{h}{6}(K_1+2K_2+2K_3+K_4),\\
K_1=y_n-\dfrac{2x_n}{y_n},\\
K_2=y_n+\dfrac{h}{2}K_1-\dfrac{2x_n+h}{y_n+\dfrac{h}{2}K_1},\\
K_3=y_n+\dfrac{h}{2}K_2-\dfrac{2x_n+h}{y_n+\dfrac{h}{2}K_2},\\
K_4=y_n+hK_3-\dfrac{2(x_n+h)}{y_n+hK_3}.
\end{cases}
\]

表 5.4 列出计算结果 \(y_n\)，其中 \(y(x_n)\) 仍表示准确解。

\href{../../source-images/pdf-132.jpeg}{原页}

比较例 5.4 和例 5.2 的计算结果，显然以四阶 Runge-Kutta 方法的精度为高。注意，虽然四阶 Runge-Kutta 方法的计算量（每一步要四次计算函数 \(f\)）比改进的 Euler 方法（它是一种二阶 Runge-Kutta 方法，每一步只要两次计算函数 \(f\)）大一倍，但由于这里放大了步长 \((h=0.2)\)，造出表 5.4 和表 5.2 所耗费的计算量几乎相同。这个例子又一次显示了选择算法的重要意义。

\textbf{表 5.4}

{\def\LTcaptype{none} % do not increment counter
\begin{longtable}[]{@{}lll@{}}
\toprule\noalign{}
\(x_n\) & \(y_n\) & \(y(x_n)\) \\
\midrule\noalign{}
\endhead
\bottomrule\noalign{}
\endlastfoot
0.2 & 1.1832 & 1.1832 \\
0.4 & 1.3417 & 1.3416 \\
0.6 & 1.4833 & 1.4832 \\
0.8 & 1.6125 & 1.6125 \\
1.0 & 1.7321 & 1.7321 \\
\end{longtable}
}

然而值得指出的是，Runge-Kutta 方法的推导基于 Taylor 展开方法，因而它要求所求的解具有较好的光滑性质。反之，如果解的光滑性差，那么，使用四阶 Runge-Kutta 方法可能反而不如使用改进的 Euler 方法求得的数值解的精度高。实际计算时，应当针对问题的具体特点选择合适的算法。

\subsubsection{5.3.6 变步长的 Runge-Kutta 方法}\label{ux53d8ux6b65ux957fux7684-runge-kutta-ux65b9ux6cd5}

单从每一步看，步长越小，截断误差就越小；但随着步长的缩小，在一定求解范围内所要完成的步数就增加了。步数的增加不但引起计算量的增大，而且可能导致舍入误差的严重积累。因此同积分的数值计算一样，微分方程的数值解法也有个选择步长的问题。

在选择步长时，需要考虑两个问题：

（1）怎样衡量和检验计算结果的精度？

（2）如何依据所获得的精度处理步长？

考察经典的四阶 Runge-Kutta 格式 (5.3.13)。从节点 \(x_n\) 出发，先以 \(h\) 为步长求出一个近似值，作为 \(y_{n+1}^{(h)}\)，由于格式的局部截断误差为 \(O(h^5)\)，故有

\[
y(x_{n+1})-y_{n+1}^{(h)}\approx Ch^5;\tag{5.3.14}
\]

然后将步长折半，即取 \(\dfrac{h}{2}\) 为步长从 \(x_n\) 跨两步到 \(x_{n+1}\)，再求得一个近似值 \(y_{n+1}^{(h/2)}\)，每跨一步的截断误差是 \(C\left(\dfrac{h}{2}\right)^5\)，因此有

\[
y(x_{n+1})-y_{n+1}^{(h/2)}=2C\left(\frac{h}{2}\right)^5;\tag{5.3.15}
\]

比较式 (5.3.14) 和式 (5.3.15) 可以看到，步长折半后，误差大约减小到 \(\dfrac{1}{16}\)，即有

\[
\frac{y(x_{n+1})-y_{n+1}^{(h/2)}}{y(x_{n+1})-y_{n+1}^{(h)}}\approx\frac{1}{16}.
\]

由此易得下列事后估计式

\[
y(x_{n+1})-y_{n+1}^{(h/2)}\approx\frac{1}{15}(y_{n+1}^{(h/2)}-y_{n+1}^{(h)}).
\]

\href{../../source-images/pdf-133.jpeg}{原页}

这样，可以通过检查步长折半前后两次计算结果的偏差

\[
\Delta=|y_{n+1}^{(h/2)}-y_{n+1}^{(h)}|
\]

来判定所选的步长是否合适。具体地说，将分以下两种情况处理：

（1）对于给定的精度 \(\varepsilon\)，如果 \(\Delta>\varepsilon\)，应反复将步长折半进行计算，直至 \(\Delta<\varepsilon\) 为止，这时取最终得到的 \(y_{n+1}^{(h/2)}\) 作为结果；

（2）如果 \(\Delta<\varepsilon\)，则反复将步长加倍，直至 \(\Delta>\varepsilon\) 为止，这时再将步长折半一次，就得到所要的结果。

这种通过加倍或折半处理步长的方法称为\textbf{变步长方法}。表面上看，为了选择步长，每一步的计算量增加了，但从总体考虑往往是合算的。

\subsection{本节覆盖原页的校注记录}\label{ux672cux8282ux8986ux76d6ux539fux9875ux7684ux6821ux6ce8ux8bb0ux5f55}

下列记录按原页关联，含同页相邻小节的校注；计算时同时读取上文校注和完整原页。

\begin{itemize}
\tightlist
\item
  \texttt{ch05a-E003}：原书 PDF 页 129
\end{itemize}

\end{document}
